В современном мире цифровых технологий компьютерные игры стали неотъемлемой частью глобальной культуры и индустрии развлечений, превратившись из нишевого хобби в мощный социальный и экономический феномен. По данным аналитиков, мировой игровой рынок продолжает демонстрировать устойчивый рост, а игры сегодня "--- это сложные произведения, сочетающие в себе элементы искусства, повествования, инженерии и психологии.

В этом контексте процесс разработки игры предстает как комплексная творческо"=техническая задача, требующая слаженной междисциплинарной работы. Он включает в себя концептуализацию идеи, проектирование геймплея и нарратива, создание художественного и звукового оформления, программирование игровых систем, а также тестирование и продвижение продукта. Каждый этап этого процесса критически важен для достижения конечной цели "--- создания увлекательного, полированного и запоминающегося игрового опыта, который найдет отклик у целевой аудитории.

Проект посвящен созданию компьютерной игры"=головоломки <<Одноплатник>>, которая создана для изучения будущими студентами факультета специальности <<Информатика и вычислительная техника>>. В игре реализованы элементы  головоломки с продуманной визуальной составляющей, написанной на языке программирования С++.

Сегодня компьютерные игры "--- не только развлечение. Они становятся мощным инструментом обучения, тренажёрами для инженеров, симуляторами сложных систем. Особенно актуален этот тренд в техническом образовании, где абстрактные концепции "--- архитектура компьютера, интерфейсы ввода-вывода или принципы работы периферии "--- зачастую трудны для восприятия на первых курсах.

В этой связи обучающие игры, приобретают особое значение. Они позволяют студенту не просто запомнить, что HDMI используется для передачи видео, а испытать этот принцип на практике. Проект <<Одноплатник>> представляет собой такую среду "--- минималистичную, но содержательную, где каждый провод, каждый порт и каждое соединение несут учебную нагрузку.

Цель проекта "--- разработка образовательной видеоигры"=головоломки, способствующей формированию интуитивного понимания у первокурсников базовых принципов работы вычислительных систем и периферийных устройств через моделирование реальных задач.